\section{A fixed parameter in the conic chart}

For a nonzero rational pole input $b$, write $B=17b$ and
\[
 e=\frac{\tau^2-B}{\tau^2+3B},\qquad
 p=\frac{-4B\tau}{\tau^2+3B}.
\]
Fixing $\tau$ is different from fixing $e$, since $e$ still varies with $b$.
The following exact example rules out the choice $\tau=1$, as well as
$\tau=-1$, as a complete specialization on all prepared pole inputs.

\begin{proposition}
The rational number
\[
 b=\left(\frac{204}{61}\right)^2=\frac{41616}{3721}
\]
has square reduced denominator, satisfies $0<b<17$, is a square in
$\mathbf Q_2,\mathbf Q_3,\mathbf Q_{17}$, and has respective valuations
$4,2,2$. For $\tau=1$ or $\tau=-1$, the equation
\[
 Bq^2=(s^2-Be(B-b^2))(s^2+3Be(B-b^2))
\]
has no solution $(s,q)\in\mathbf Q^2$.
\end{proposition}
\begin{proof}
The preparation assertions follow directly from $b=(204/61)^2$.
For either choice of $\tau$,
\[
 e=-\frac{703751}{2126137}.
\]
Put $M=Be(B-b^2)$ and $N=BM$. Direct rational arithmetic gives
\[
 N=nk^2,\qquad
 n=-32380801574194567,\qquad
 k=\frac{144324288}{29438154846217}.
\]
A hypothetical solution with $s\ne0$ would give the point
\[
 x=\frac{B}{k^2}s^2,\qquad y=\frac{B^2}{k^3}sq
\]
on
\[
 E_n:\quad y^2=x(x-n)(x+3n).
\]
Its nonzero abscissa would belong to the squareclass $17$.

Exact two-descent gives $\operatorname{rank}E_n(\mathbf Q)=0$.
The accompanying reproducible PARI calculation returns
\texttt{ellrank} equal to \texttt{[0,0,0,[]]}; it uses no analytic-rank
or Birch--Swinnerton-Dyer hypothesis. Independent counts at the good
primes $5,7,11,23,29,31$ are respectively
\[
                         4,8,8,32,24,24.
\]
Their greatest common divisor is four. Together with the visible full
two-torsion this proves
\[
 E_n(\mathbf Q)=\{O,(0,0),(n,0),(-3n,0)\}.
\]
Neither nonzero torsion abscissa has squareclass $17$: $n$ is negative,
and $-3n/17$ is not a rational square. This excludes $s\ne0$.
If $s=0$, the original equation becomes $Bq^2=-3M^2<0$, which is
impossible over $\mathbf R$. This proves the assertion.
\end{proof}

This proposition concerns all prepared inputs. It does not decide
whether a fixed conic parameter can be made complete on a more
restricted family chosen simultaneously with the comparison primes.
